%% Jaren Untuña --- Curriculum Vitae (English)
%% Compile:  tectonic cv-en.tex
%% Plantilla propia: una columna, sin iconos, sin barras de "nivel de skill".
%% Está pensada para pasar un ATS y para que una persona lo lea en 20 segundos.

\documentclass[10.5pt,a4paper]{article}

\usepackage[T1]{fontenc}
\usepackage{tgheros}
\renewcommand*\familydefault{\sfdefault}
\usepackage[utf8]{inputenc}
\usepackage[spanish,es-noshorthands]{babel}
\usepackage[
  a4paper,
  top=1.5cm, bottom=1.5cm, left=1.6cm, right=1.6cm,
  footskip=0.8cm
]{geometry}
\usepackage{titlesec}
\usepackage{enumitem}
\usepackage{xcolor}
\usepackage{hyperref}
\usepackage{tabularx}
\usepackage{microtype}

\definecolor{ink}{HTML}{111318}
\definecolor{muted}{HTML}{55606B}
\definecolor{rule}{HTML}{C9D0D8}
\definecolor{accent}{HTML}{0E7C6B}

\hypersetup{
  colorlinks=true,
  urlcolor=accent,
  linkcolor=accent,
  pdftitle={Johan Untuna - Ingeniero de Software - CV},
  pdfauthor={Johan Untuna},
  pdfsubject={Curriculum Vitae},
  pdfkeywords={software engineer, python, django, data pipelines, api integration, security}
}

\color{ink}
\pagestyle{empty}
\setlength{\parindent}{0pt}
\setlength{\emergencystretch}{2em}

%% ---------------------------------------------------------------- headings
\titleformat{\section}
  {\normalfont\large\bfseries\scshape}
  {}{0pt}{}
  [\vspace{-0.55em}{\color{rule}\rule{\linewidth}{0.5pt}}]
\titlespacing*{\section}{0pt}{1.15em}{0.6em}

%% Entry header: title on the left, dates flush right.
\newcommand{\entry}[2]{%
  \begin{tabularx}{\linewidth}{@{}Xr@{}}
    \textbf{#1} & {\small\color{muted}#2}
  \end{tabularx}%
  \vspace{0.15em}
}

%% One-line subtitle under an entry header.
\newcommand{\subentry}[1]{{\small\color{muted}#1}\vspace{0.25em}}

\newenvironment{points}
  {\begin{itemize}[leftmargin=1.1em, itemsep=0.18em, topsep=0.25em, parsep=0pt, label=\textcolor{accent}{\small$\bullet$}]}
  {\end{itemize}}

%% Skill row: bold category, then the list.
\newcommand{\skillrow}[2]{%
  \begin{tabularx}{\linewidth}{@{}p{3.5cm}X@{}}
    \textbf{#1} & #2
  \end{tabularx}%
  \vspace{0.15em}
}

\begin{document}

%% ------------------------------------------------------------------ header
{\Huge\bfseries Johan Untuña}\\[0.35em]
{\large\color{muted} Ingeniero de Software --- Pipelines de Datos, Integración de APIs e IA Aplicada}\\[0.7em]
{\small
Latacunga, Ecuador (GMT-5) \textbullet{} Disponible en remoto \\
\href{mailto:jaren22uv@gmail.com}{jaren22uv@gmail.com} \textbullet{}
\href{https://wa.me/593963607760}{+593 96 360 7760} \\
\href{https://github.com/JohanUV}{github.com/JohanUV} \textbullet{}
\href{https://www.linkedin.com/in/jaren-untu\%C3\%B1a-valiente-762799211/}{LinkedIn} \textbullet{}
\href{https://tryhackme.com/p/jaren22uv}{TryHackMe}
}

\vspace{0.4em}
{\color{rule}\rule{\linewidth}{0.9pt}}

%% ----------------------------------------------------------------- summary
\section*{Perfil}
Estudiante de Ingeniería de Software (5.º semestre, ESPE Latacunga) que construye y entrega
sistemas completos, no ejercicios. Mi trabajo se centra en extraer datos de fuentes que se
resisten a ser leídas ---APIs estatales sin documentar, portales detrás de captchas, proveedores
con esquemas incompatibles--- y volver el resultado lo bastante confiable como para decidir con
él. Diseño con control de acceso, auditoría y protección de datos desde el inicio. Dos de mis
proyectos fueron encargos de negocios locales; el resto es público en GitHub con la
arquitectura documentada. En transición hacia seguridad ofensiva y OSINT.

%% ---------------------------------------------------------------- projects
\section*{Proyectos Destacados}

\entry{Verum --- Plataforma de Debida Diligencia para Contratación}{2026}
\subentry{Python, Flask, React, REST, RBAC \textbullet{} \href{https://github.com/JohanUV/Verum-EC}{github.com/JohanUV/Verum-EC}}
\begin{points}
  \item Integré \textbf{nueve fuentes públicas ecuatorianas} ---Función Judicial, SRI, SUPA,
        ANT y listas de sanción OFAC/ONU--- en un informe de riesgo consolidado consultado
        por cédula.
  \item Resolví las fuentes JSF protegidas con captcha mediante \textbf{relay humano} en vez de
        romperlas: la app muestra el captcha al usuario y continúa la sesión autenticada,
        manteniendo la integración dentro de lo que cada fuente permite.
  \item Etiqueté cada fuente con un estado de fiabilidad explícito (API en vivo, relay de
        captcha, asistida) para que el informe nunca exagere su propia cobertura.
  \item Construida para la LOPDP: permisos por rol, historial de auditoría completo de quién
        consultó a quién, y exportación a PDF con folio y marca de agua para trazabilidad.
\end{points}

\entry{Job Radar --- Pipeline Automatizada de Inteligencia de Empleo}{2026}
\subentry{Django, DRF, PostgreSQL, n8n, React, Docker, Gemini \textbullet{} \href{https://github.com/JohanUV/Job-Radar}{github.com/JohanUV/Job-Radar}}
\begin{points}
  \item Pipeline programada que recolecta vacantes de tres APIs de empleo cada seis horas,
        normalizando cada proveedor a un esquema común mediante conectores por fuente.
  \item Hice los duplicados \textbf{estructuralmente imposibles} forzando un hash SHA-256 único
        de la URL a nivel de base de datos: reejecutar la pipeline no puede corromper datos,
        sin importar la concurrencia.
  \item Separé la orquestación (n8n: programación, descarga en paralelo, reintentos) de las
        reglas de negocio (Django: validación, deduplicación, persistencia) --- añadir una fuente
        no toca el backend.
  \item Añadí una etapa LLM que puntúa afinidad CV--vacante 0--100 con Gemini, pull-based e
        idempotente: puede caerse y reanudar sin doble facturación del modelo; guarda el
        razonamiento y la versión del modelo junto a cada puntaje.
\end{points}

\entry{Análisis de Procesos Judiciales --- API Estatal por Ingeniería Inversa}{2026}
\subentry{Python, HTTP traffic analysis, API specification \textbullet{} \href{https://github.com/JohanUV/Sistema-Analisis-Judicial}{github.com/JohanUV/Sistema-Analisis-Judicial}}
\begin{points}
  \item Reconstruí el contrato no documentado del servicio detrás del portal judicial del
        Ecuador a partir del tráfico observado: dos endpoints, forma completa del payload,
        paginación y los headers que el servidor valida realmente.
  \item Identifiqué que el servicio valida \textbf{procedencia} (Origin/Referer) y no identidad,
        y que su campo de captcha se acepta como constante --- un hallazgo de confianza en el
        cliente documentado en la capa de servicio.
  \item Documenté el resultado como referencia de API reutilizable en vez de dejarlo implícito
        en el código de un scraper.
\end{points}

\entry{Class Attention --- Detección de Atención en el Aula}{2026}
\subentry{Python, machine learning, computer vision, privacy by design}
\begin{points}
  \item Sistema de ML que estima la atención de los estudiantes a partir de video, validado en
        sesiones reales de clase, dando al docente una señal accionable a mitad de lección.
  \item Diseñado para que la salida útil sea un \textbf{agregado del aula}, nunca un registro por
        estudiante --- el riesgo de privacidad se elimina en la arquitectura, no por política.
  \item La protección de identidades fue una restricción desde la primera versión: los sujetos
        son menores y una filtración sería irreparable.
\end{points}

%% ------------------------------------------------------------ client work
\section*{Trabajo para Clientes}

\entry{Cottullari --- Operadora de Turismo, Latacunga}{2026}
\subentry{React, React Router, Vite, PHP, MySQL \textbullet{} \href{https://github.com/JohanUV/cottullari-react}{github.com/JohanUV/cottullari-react}}
\begin{points}
  \item Construí una SPA de siete rutas para una empresa de turismo: flota, servicios, destinos
        y cotizaciones, con un formulario que persiste las solicitudes en MySQL.
  \item Resolví las cinco páginas de destino con una sola ruta parametrizada: añadir un destino
        es dato, no un componente nuevo.
\end{points}

\entry{Mashka Box --- Box de CrossFit, Latacunga}{2026}
\subentry{HTML, CSS, responsive \textbullet{} \href{https://github.com/JohanUV/mashka-box}{github.com/JohanUV/mashka-box}}
\begin{points}
  \item Landing page mobile-first hecha para convertir una visita en una clase de prueba; HTML
        y CSS a mano, sin framework, porque una página sin estado de aplicación no debería
        hacer que el visitante descargue uno.
\end{points}

%% ------------------------------------------------------------------ skills
\section*{Competencias Técnicas}
\skillrow{Lenguajes}{Python, TypeScript, JavaScript, SQL, HTML, CSS}
\skillrow{Backend}{Django, Django REST Framework, Flask, REST API design}
\skillrow{Frontend}{React, React Router, Vite, Astro, Tailwind CSS}
\skillrow{Datos}{PostgreSQL, SQLite, Drizzle ORM, Pandas, web scraping, normalización de esquemas}
\skillrow{IA y ML}{Gemini API, OpenAI API, Ollama, OpenCV, scikit-learn, diseño de prompts}
\skillrow{Automatización}{n8n, pipelines programadas, arquitecturas de conectores, bots de Telegram}
\skillrow{Infraestructura}{Linux, Bash, Docker, Docker Compose, Terraform, LocalStack, AWS Lambda, Vercel, Git}
\skillrow{Seguridad}{Control de acceso (RBAC), auditoría, autenticación por API key, cumplimiento de protección de datos}
\skillrow{Aprendiendo}{Seguridad ofensiva, Burp Suite, Nmap, protocolos de red, OSINT}

%% --------------------------------------------------------------- education
\section*{Formación}
\entry{Universidad de las Fuerzas Armadas ESPE --- Latacunga}{2024 --- 2028 (estimado)}
\subentry{Ingeniería de Software \textbullet{} Cursando 5.º semestre}

%% --------------------------------------------------------------- languages
\section*{Idiomas}
\skillrow{Español}{Nativo}
\skillrow{Inglés}{B2 --- competencia profesional, cómodo en conversación técnica}

\end{document}
